\section{The Aldous--Lyons conjecture: subgroup tests}
\label{ch:subgroup-tests}

This chapter is the plan and the progress of Phase~6 of \texttt{planning/aldous-lyons-track.md}:
the companion paper~\cite{BCLV24} of Bowen, Chapman, Lubotzky and Vidick, which with
Chapter~\ref{ch:tailored} refutes the Aldous--Lyons conjecture. Line numbers I:$n$ refer to the
LaTeX source of its arXiv version, which is the authority on the mathematics. Its argument has
two halves. Main Theorem~II (Theorem~I:2126), finitary, associates with a tailored game a
subgroup test whose sofic value is $1$ when the game has a perfect ZPC strategy and bounded away
from $1$ when the game's value is at most $\frac12$. Main Theorem~I (I:594) shows that under the
conjecture the sofic value is approximable by a computable function. Together with a halting
reduction to tailored games they contradict the undecidability of the halting problem.

\emph{Status.} The definitions are formalized (\texttt{MIPRE/SubgroupTestValue.lean},
self-contained, and \texttt{MIPRE/Tailored/Sofic/}): subgroup tests and their sofic value, the
Aldous--Lyons conjecture, and the associated test of a tailored game. Proved: the associated
test is computable from the game (Lemma~\ref{lem:assoc-test-primrec}), and Main Theorem~II with a
halting reduction makes the sofic value not approximable
(Theorem~\ref{thm:sofic-not-approximable}); and the perturbation half of Main Theorem~II's
soundness (Proposition~\ref{prop:assoc-perturb}); Main Theorem~II itself
(Theorem~\ref{thm:main-theorem-ii}); and so the sofic value is not approximable from a halting
reduction to tailored games (Corollary~\ref{cor:sofic-not-approximable-ar}). Main Theorem~I
(Theorems~\ref{thm:main-theorem-i-lower} and~\ref{thm:main-theorem-i-upper}) is formalized too, so
the Aldous--Lyons conjecture is false from a halting reduction to tailored games, and from a
tailored answer reduction (Corollary~\ref{cor:aldous-lyons-false-ar}). Paper~II's main theorem is
proved (Theorem~\ref{thm:tmip-re}), so the conjecture is false
(Corollary~\ref{cor:aldous-lyons-false-final}).

\subsection{Subgroup tests and the conjecture}

\begin{definition}[Subgroup tests and their sofic value; I:439--520]
\label{def:subgroup-test}
\lean{SubgroupTestValue.SubgroupTestData, SubgroupTestValue.SubgroupTestData.equivTuple,
  SubgroupTestValue.FiniteAction, SubgroupTestValue.FiniteAction.letterPerm,
  SubgroupTestValue.FiniteAction.wordPerm, SubgroupTestValue.FiniteAction.InStab,
  SubgroupTestValue.FiniteAction.LitHolds, SubgroupTestValue.FiniteAction.Passes,
  SubgroupTestValue.FiniteAction.passProb, SubgroupTestValue.SubgroupTestData.totalWeight,
  SubgroupTestValue.SubgroupTestData.value, SubgroupTestValue.SubgroupTestData.valSof}
\leanok
A subgroup test over the free group $F$ on $n$ generators is a finite list of weighted
challenges, each a finite list $K$ of words and a decision on the subsets of $K$, here in
disjunctive normal form over the literals ``the $i$-th word of $K$ lies in (or out of) the
subgroup''. A finite action $\sigma$ of $F$ on a set of size $N \geq 1$ gives the finitely
described IRS $\Phi(\sigma)$, the stabilizer of a uniform point; its value against the test is the
probability that the stabilizer passes a challenge drawn with probability proportional to its
weight. The sofic value of the test is the supremum of the values of finite actions.
\end{definition}

\begin{definition}[Approximating the sofic value]
\label{def:sofic-approximable}
\uses{def:subgroup-test}
\lean{SubgroupTestValue.SofValueApproximable}
\leanok
The sofic value is approximable when a computable function $f$ of a test $T$ and $k \in \N$
satisfies $|\mathrm{val}_{\mathrm{sof}}(T) - f(T, k)/(k+1)| \leq 1/(k+1)$ for all $T$ and $k$.
\end{definition}

\begin{definition}[Invariant random subgroups and the Aldous--Lyons conjecture; I:467--530]
\label{def:aldous-lyons}
\uses{def:subgroup-test}
\lean{SubgroupTestValue.IsSubgroupIndicator, SubgroupTestValue.SubgroupSpace,
  SubgroupTestValue.conj, SubgroupTestValue.IRS, SubgroupTestValue.stab,
  SubgroupTestValue.finDescIRS, SubgroupTestValue.AldousLyons}
\leanok
The space $\mathrm{Sub}(F)$ of subgroups of the free group on $s$ generators is a closed subspace
of $\{0,1\}^F$ with the product topology. An invariant random subgroup is a Borel probability
measure on it invariant under conjugation; the finitely described ones are the measures
$\Phi(\sigma)$ of finite actions. The Aldous--Lyons conjecture asserts that for every $s$ the
finitely described IRSs are weak-$*$ dense in the IRSs.
\end{definition}

\subsection{The test associated with a tailored game}

\begin{definition}[The associated subgroup test; I:2086]
\label{def:assoc-test}
\uses{def:subgroup-test,def:tailored-game-data}
\lean{MIPRE.Tailored.Sofic.invW, MIPRE.Tailored.Sofic.commW, MIPRE.Tailored.Sofic.genW,
  MIPRE.Tailored.Sofic.allBits, MIPRE.Tailored.Sofic.genJ, MIPRE.Tailored.Sofic.wts,
  MIPRE.Tailored.Sofic.wts_eq, MIPRE.Tailored.Sofic.genX, MIPRE.Tailored.Sofic.nGen,
  MIPRE.Tailored.Sofic.wJ, MIPRE.Tailored.Sofic.wX, MIPRE.Tailored.Sofic.varsAt,
  MIPRE.Tailored.Sofic.fixedLits, MIPRE.Tailored.Sofic.readVars,
  MIPRE.Tailored.Sofic.readWords, MIPRE.Tailored.Sofic.consAt, MIPRE.Tailored.Sofic.consWord,
  MIPRE.Tailored.Sofic.consWords, MIPRE.Tailored.Sofic.words, MIPRE.Tailored.Sofic.clause,
  MIPRE.Tailored.Sofic.clauses, MIPRE.Tailored.Sofic.challenge, MIPRE.Tailored.Sofic.assocTest,
  MIPRE.Tailored.Sofic.genPerm, MIPRE.Tailored.Sofic.Checks}
\leanok
The test associated with a tailored game has a generator $J$ and one generator for each variable
of each vertex, and one challenge per question pair, with the pair's weight. A subgroup $H$
passes the challenge at $(x, y)$ when: $J \notin H$, $J^2 \in H$ and $J$ commutes in $H$ with the
variables at $x$ and $y$ (Check~1); the variables are involutions and those at a vertex commute
(Check~2); each readable variable $X$ has $X \in H$ or $J X \in H$, which reads off a value $r$
of the readable variables (Check~3); and the word $J^{c_J} \prod X^{c}$ of every constraint
$c \in L_{xy}(r)$ lies in $H$ (Check~4). It is written as one clause per readable value $r$, and
the constraint words are deduplicated, which bounds their number by $2^{2\Lambda + 1} + 1$. An
action passes Checks~1--3 everywhere when $J$ is a fixed-point-free involution commuting with the
variables, these are involutions commuting at each vertex, and each readable variable acts at
every point as the identity or as $J$.
\end{definition}

\begin{lemma}[The associated test is computable from the game]
\label{lem:assoc-test-primrec}
\uses{def:assoc-test,lem:tailored-conversion}
\lean{MIPRE.Tailored.Sofic.allBits_eq_iterate, MIPRE.Tailored.Sofic.dedup_eq_foldr,
  MIPRE.Tailored.Sofic.zipIdx_map_eq, MIPRE.Tailored.Sofic.primrec_consWord,
  MIPRE.Tailored.Sofic.primrec_consWords, MIPRE.Tailored.Sofic.primrec_clause,
  MIPRE.Tailored.Sofic.primrec_clauses, MIPRE.Tailored.Sofic.totalWeight_eq,
  MIPRE.Tailored.Sofic.primrec_assocTest}
\leanok
The map from a tailored game description to the description of its associated test is
primitive recursive.
\end{lemma}

\begin{proof}
\leanok
Each construction is a composition of list operations. Three are rewritten first: the readable
values as an iterate, deduplication as a fold, and the enumeration of the fixed literals by a
range.
\end{proof}

\begin{lemma}[Hamming distance of permutations]
\label{lem:perm-hamming}
\lean{MIPRE.Tailored.Sofic.diffCard, MIPRE.Tailored.Sofic.dH, MIPRE.Tailored.Sofic.diffCard_triangle,
  MIPRE.Tailored.Sofic.diffCard_mul, MIPRE.Tailored.Sofic.diffCard_inv,
  MIPRE.Tailored.Sofic.card_commutator_moves, MIPRE.Tailored.Sofic.dH_triangle,
  MIPRE.Tailored.Sofic.dH_mul, MIPRE.Tailored.Sofic.dH_inv}
\leanok
The normalized Hamming distance of two permutations of a finite set, the fraction of points
where they differ, satisfies the triangle inequality, is subadditive on products, and is
invariant under inversion.
\end{lemma}

\begin{proof}
\leanok
Counting points.
\end{proof}

\begin{lemma}[Stability of involutions, fixed-point-free involutions and commuting families;
I:2561--2649]
\label{lem:perm-stability}
\uses{lem:perm-hamming}
\lean{MIPRE.Tailored.Sofic.invFix, MIPRE.Tailored.Sofic.invFix_invol,
  MIPRE.Tailored.Sofic.diffCard_invFix, MIPRE.Tailored.Sofic.fpfFix,
  MIPRE.Tailored.Sofic.fpfFix_invol, MIPRE.Tailored.Sofic.fpfFix_free,
  MIPRE.Tailored.Sofic.diffCard_fpfFix, MIPRE.Tailored.Sofic.ordProd,
  MIPRE.Tailored.Sofic.commFix, MIPRE.Tailored.Sofic.commFix_invol,
  MIPRE.Tailored.Sofic.commFix_comm, MIPRE.Tailored.Sofic.commFix_J,
  MIPRE.Tailored.Sofic.commFix_readable, MIPRE.Tailored.Sofic.diffCard_commFix}
\leanok
A permutation $\pi$ is within $d(\pi^2, 1)$ of an involution; an involution of an even set almost
without fixed points is close to one without, given a fixed-point-free involution commuting with
it; and a family of involutions at a vertex, together with $J$, can be made into commuting
involutions commuting with $J$, the readable ones acting as the identity or as $J$, by changing
them on at most $2^{\ell + 1}$ times the points at which some relation fails.
\end{lemma}

\begin{proof}
\leanok
The first two as in the paper. The third replaces Claim~I:2600 and Glebsky--Rivera
(Claim~I:2631): restrict the variables to the set of points from which every ordered product of
them reaches only points where all relations hold, invariant under the variables and $J$, and
let them act as the identity elsewhere.
\end{proof}

\begin{lemma}[Robustness of the value; Claim~I:1480]
\label{lem:test-robust}
\uses{def:subgroup-test,lem:perm-hamming}
\lean{MIPRE.Tailored.Sofic.wordCount, MIPRE.Tailored.Sofic.sig, MIPRE.Tailored.Sofic.withPerms,
  MIPRE.Tailored.Sofic.diffCard_wordPerm, MIPRE.Tailored.Sofic.abs_passProb_sub_le,
  MIPRE.Tailored.Sofic.abs_value_sub_le}
\leanok
For two actions of the generators on the same finite set,
$|\mathrm{val}(T, \sigma) - \mathrm{val}(T, \sigma')| \leq \sum_k S_T(k)\, d(\sigma_k,
\sigma'_k)$, where the significance $S_T(k)$ is the expected number of occurrences of the
generator $k$ in the words of a challenge.
\end{lemma}

\begin{proof}
\leanok
A word moves a point differently under the two actions only if one of its letters does at an
intermediate point; letters and their inverses cost the same.
\end{proof}

\begin{lemma}[Significance of the associated test; Proposition~I:2283]
\label{lem:assoc-significance}
\uses{def:assoc-test,lem:test-robust}
\lean{MIPRE.Tailored.Sofic.sigBound, MIPRE.Tailored.Sofic.sig_le, MIPRE.Tailored.Sofic.mdeg,
  MIPRE.Tailored.Sofic.sum_mdeg, MIPRE.Tailored.Sofic.sig_genX_le}
\leanok
In the associated test every generator has significance at most a bound $S(\Lambda)$ of order
$\Lambda 4^{\Lambda}$, and the variables at a vertex $x$ at most $S(\Lambda)$ times the weight of
the question pairs containing $x$.
\end{lemma}

\begin{proof}
\leanok
The words of the challenge at $(x, y)$ use only $J$ and the variables at $x$ and $y$; their
number and lengths are bounded, the constraint words by deduplication.
\end{proof}

\begin{proposition}[Perturbing to pass Checks~1--3; Proposition~I:2267]
\label{prop:assoc-perturb}
\uses{lem:perm-stability,lem:test-robust,lem:assoc-significance}
\lean{MIPRE.Tailored.Sofic.double, MIPRE.Tailored.Sofic.value_double, MIPRE.Tailored.Sofic.dblSwap,
  MIPRE.Tailored.Sofic.qw, MIPRE.Tailored.Sofic.qwTot, MIPRE.Tailored.Sofic.qwTot_pos,
  MIPRE.Tailored.Sofic.value_eq, MIPRE.Tailored.Sofic.sig_eq,
  MIPRE.Tailored.Sofic.fixed_of_passes, MIPRE.Tailored.Sofic.read_of_passes,
  MIPRE.Tailored.Sofic.npass, MIPRE.Tailored.Sofic.SwapData, MIPRE.Tailored.Sofic.J₂,
  MIPRE.Tailored.Sofic.Xfix, MIPRE.Tailored.Sofic.perturbed,
  MIPRE.Tailored.Sofic.checks_perturbed, MIPRE.Tailored.Sofic.card_bad_le,
  MIPRE.Tailored.Sofic.lose, MIPRE.Tailored.Sofic.loseAt, MIPRE.Tailored.Sofic.weighted_dist_le,
  MIPRE.Tailored.Sofic.value_perturbed_ge, MIPRE.Tailored.Sofic.numeric_bound,
  MIPRE.Tailored.Sofic.Cp, MIPRE.Tailored.Sofic.exists_checks_of_value}
\leanok
If an action has value at least $1 - \varepsilon$ against the associated test of a tailored game
with answer length $\Lambda$, then an action passing Checks~1--3 everywhere has value at least
$1 - 300\,(\Lambda + 1)^4\, 2^{4(\Lambda + 1)}\, \varepsilon$.
\end{proposition}

\begin{proof}
\leanok
Double the action first, so that the set is even and carries a fixed-point-free involution
commuting with every generator. Make $J$ a fixed-point-free involution, then repair the variables
vertex by vertex (Lemma~\ref{lem:perm-stability}), bounding the points changed at a vertex by its
losing probability, and conclude by robustness and significance.
\end{proof}

\begin{definition}[Z-aligned signed-permutation strategies without commutation along edges]
\label{def:zstrat}
\uses{def:tailored-game-data,def:sync-strategy,def:sync-value,def:signed-perm,lem:signed-perm,def:fourier-forms,lem:fourier-pvm,lem:tailored-bridge}
\lean{MIPRE.Tailored.Sofic.ZStrat, MIPRE.Tailored.Sofic.ZStrat.proj,
  MIPRE.Tailored.Sofic.ZStrat.value, MIPRE.Tailored.Sofic.ZStrat.toSync,
  MIPRE.Tailored.Sofic.ZStrat.value_toSync, MIPRE.Tailored.Sofic.ZStrat.value_le_gameValue,
  MIPRE.Tailored.Sofic.ZStrat.rej, MIPRE.Tailored.Sofic.ZStrat.value_eq_sum_rej,
  TailoredGameValue.PermStrategy.toZStrat, TailoredGameValue.PermStrategy.value_toZStrat}
\leanok
A permutation strategy for a tailored game (signed permutation observables, involutions
commuting at each vertex, diagonal at the readable variables, the identity beyond each vertex's
length) need not commute along edges; its Fourier-transform measurements still form a
synchronous strategy, with the same value, so its value is at most the game's synchronous value.
\end{definition}

\begin{lemma}[The trace identity]
\label{lem:trace-identity}
\uses{def:zstrat,def:signed-perm,lem:signed-perm,def:fourier-forms,lem:fourier-pvm}
\lean{MIPRE.Tailored.Sofic.ZStrat.rsgn, MIPRE.Tailored.Sofic.ZStrat.sum_trace_proj_mul_mul_proj,
  MIPRE.Tailored.Sofic.consBit, MIPRE.Tailored.Sofic.satisfies_iff_consBit,
  MIPRE.Tailored.Sofic.ZStrat.consMat, MIPRE.Tailored.Sofic.ZStrat.proj_mul_consMat_mul_proj}
\leanok
For such a strategy, the sum over the answer pairs whose readable parts lie in a set $R$ of
$\mathrm{Tr}(P_a\, X\, Q_b)$ is the sum of the diagonal entries $X_{jj}$ over the basis points $j$
whose readable signs lie in $R$. For a constraint $c$ with matrix $M_c = (-1)^{c_J} U^{\alpha}
V^{\beta}$, the probability that the readable parts are $r$ and $c$ is violated is
$\frac1m \sum_{j \in R_r} \frac{1 - (M_c)_{jj}}{2}$.
\end{lemma}

\begin{proof}
\leanok
A basis point lies in the support of $P_a$ only if its readable signs are those of $a$, and the
projections sum to the identity; the observable of a constraint is $(-1)^{\langle c, (a, b, 1)
\rangle}$ on the range of $P_a \otimes Q_b$.
\end{proof}

\begin{lemma}[Passing a challenge under Checks~1--3]
\label{lem:assoc-passes}
\uses{def:assoc-test}
\lean{MIPRE.Tailored.Sofic.passes_iff, MIPRE.Tailored.Sofic.value_eq_sum,
  MIPRE.Tailored.Sofic.length_consWords_le'}
\leanok
For an action passing Checks~1--3 everywhere, a point passes the challenge at $(x, y)$ exactly
when the constraint words of the entries whose readable part is the point's readable value fix
it; and there are at most $2^{\ell(x) + \ell(y) + 1} + 1$ constraint words.
\end{lemma}

\begin{proof}
\leanok
The fixed literals and the readable literals of the point's own clause hold, and those of every
other clause fail since $J$ moves every point.
\end{proof}

\begin{lemma}[The quotient by $J$]
\label{lem:j-quotient}
\uses{def:zstrat,lem:signed-perm,lem:trace-identity,def:assoc-test,def:subgroup-test}
\lean{MIPRE.Tailored.Sofic.Reps, MIPRE.Tailored.Sofic.rep, MIPRE.Tailored.Sofic.ind,
  MIPRE.Tailored.Sofic.ind_mul, MIPRE.Tailored.Sofic.ind_J, MIPRE.Tailored.Sofic.ind_fix_iff,
  MIPRE.Tailored.Sofic.indF, MIPRE.Tailored.Sofic.quotStrat, MIPRE.Tailored.Sofic.wordSP,
  MIPRE.Tailored.Sofic.ZStrat.toMatrix_wordSP_consWord}
\leanok
Let an action pass Checks~1--3 everywhere. Choosing one point of each $J$-orbit, every
permutation commuting with $J$ induces a signed permutation of the representatives,
multiplicatively, sending $J$ to $-1$, with sign $+$ at $r$ exactly when it fixes $r$. The
variables induce a Z-aligned signed-permutation strategy, in which the matrix of a constraint is
the signed permutation induced by its word.
\end{lemma}

\begin{proof}
\leanok
A permutation commuting with $J$ permutes the $J$-orbits; the readable variables act as the
identity or as $J$, hence diagonally.
\end{proof}

\begin{proposition}[Game value of an action passing Checks~1--3; Proposition~I:2279]
\label{prop:checks-game-value}
\uses{lem:trace-identity,lem:j-quotient,lem:assoc-passes,def:assoc-test,def:subgroup-test,lem:fourier-pvm}
\lean{MIPRE.Tailored.Sofic.Cq, MIPRE.Tailored.Sofic.ZStrat.rej_le_sum_viol,
  MIPRE.Tailored.Sofic.rej_quotStrat_le, MIPRE.Tailored.Sofic.gameValue_ge_of_checks}
\leanok
An action passing Checks~1--3 everywhere with value $1 - \varepsilon$ against the associated test
gives the game a synchronous value at least $1 - 2 \cdot 2^{2(\Lambda + 1)} \varepsilon$.
\end{proposition}

\begin{proof}
\leanok
By the trace identity, the game's rejection mass at $(x, y)$ is at most the sum over the distinct
constraint words of the fraction of representatives they do not fix, hence at most the number of
constraint words, twice, times the test's failure probability at $(x, y)$; there are at most
$2^{2(\Lambda + 1)}$ of them. This replaces the paper's Fourier bases and orbit decompositions.
\end{proof}

\begin{lemma}[Completeness: lifting a perfect strategy; Theorem~I:2126 (1)]
\label{lem:assoc-completeness}
\uses{lem:trace-identity,lem:assoc-passes,def:tailored-game-data,lem:j-quotient,def:signed-perm,lem:signed-perm,def:fourier-forms,lem:fourier-pvm}
\lean{MIPRE.Tailored.Sofic.liftHom, MIPRE.Tailored.Sofic.ZStrat.toAction,
  MIPRE.Tailored.Sofic.ZStrat.checks_toAction, MIPRE.Tailored.Sofic.ZStrat.consMat_apply_self_eq_one,
  MIPRE.Tailored.Sofic.exists_value_one_of_hasPerfectZPC}
\leanok
A perfect ZPC strategy of a tailored game lifts to an action of value $1$ against the associated
test: on twice the points, $J$ swapping the copies and each variable acting by its signed
permutation.
\end{lemma}

\begin{proof}
\leanok
Checks~1--3 hold at every point. By the trace identity at value $1$, every constraint word of a
pair of positive weight fixes the points of its readable value, so every point passes every
challenge of positive weight. Commutation along edges is not used.
\end{proof}

\begin{theorem}[Main Theorem~II; Theorem~I:2126]
\label{thm:main-theorem-ii}
\uses{def:assoc-test,prop:assoc-perturb,prop:checks-game-value,lem:assoc-completeness,
  thm:sofic-not-approximable}
\lean{MIPRE.Tailored.Sofic.value_le_one, MIPRE.Tailored.Sofic.valSof_le_one,
  MIPRE.Tailored.Sofic.valSof_eq_one_of_hasPerfectZPC, MIPRE.Tailored.Sofic.CII,
  MIPRE.Tailored.Sofic.gameValue_ge_of_valSof, MIPRE.Tailored.Sofic.gapK,
  MIPRE.Tailored.Sofic.primrec_gapK, MIPRE.Tailored.Sofic.mainTheoremII}
\leanok
Let $G$ be a tailored game with answer length $\Lambda$. If $G$ has a perfect ZPC strategy, its
associated test has sofic value $1$. If its associated test has sofic value at least
$1 - \varepsilon$, then $G$ has synchronous value at least $1 - 600 (\Lambda+1)^4 2^{6(\Lambda+1)}
\varepsilon$. So the clauses of Theorem~\ref{thm:sofic-not-approximable} hold with the primitive
recursive gap $K(\Lambda) = 600 (\Lambda+1)^4 2^{6(\Lambda+1)}$.
\end{theorem}

\begin{proof}
\leanok
Completeness is Lemma~\ref{lem:assoc-completeness}, no action having value above $1$. Soundness:
for every $\delta > 0$ an action has value at least $1 - \varepsilon - \delta$; perturb it to pass
Checks~1--3 (Proposition~\ref{prop:assoc-perturb}) and apply
Proposition~\ref{prop:checks-game-value}; then let $\delta \to 0$.
\end{proof}

\begin{corollary}[The sofic value is not approximable, from tailored answer reduction]
\label{cor:sofic-not-approximable-ar}
\uses{thm:main-theorem-ii,thm:sofic-not-approximable,cor:tmip-re-from-ar}
\lean{MIPRE.Tailored.Sofic.not_sofValueApproximable,
  MIPRE.Tailored.not_sofValueApproximable_of_answerReduction}
\leanok
A halting reduction to tailored games makes the sofic value of subgroup tests not approximable;
in particular a tailored answer reduction at input level~$5$ does.
\end{corollary}

\begin{proof}
\leanok
Theorem~\ref{thm:sofic-not-approximable} with Theorem~\ref{thm:main-theorem-ii}, and
Corollary~\ref{cor:tmip-re-from-ar} for the halting reduction.
\end{proof}

\subsection{The sofic value is not approximable}

\begin{theorem}[The sofic value is not approximable; Corollary~I:2144, measure-free]
\label{thm:sofic-not-approximable}
\uses{def:sofic-approximable,lem:assoc-test-primrec,def:tailored-halting-reduction,
  thm:halting-undecidable,lem:tailored-conversion}
\lean{MIPRE.Tailored.Sofic.MainTheoremII, MIPRE.Tailored.Sofic.le_iff_of_gap,
  MIPRE.Tailored.Sofic.not_sofValueApproximable_of}
\leanok
Given a computable halting reduction to tailored games, and the two clauses of Main Theorem~II
with a primitive recursive gap function $K(\Lambda)$, the sofic value is not approximable.
\end{theorem}

\begin{proof}
\leanok
Approximate the sofic value of the associated test of the game of a machine at $k = 4K(\Lambda)$.
A halting machine gives sofic value $1$, so the approximation is at least $k/(k+1)$; a
non-halting one gives sofic value at most $1 - 1/(2K)$, and the approximation is less than
$k/(k+1)$. The comparison is computable, and decides the halting problem.
\end{proof}

\subsection{The measure side: Main Theorem~I}

\begin{definition}[The value of a measure, and the ergodic value; I:486--499]
\label{def:erg-value}
\uses{def:subgroup-test,def:aldous-lyons}
\lean{MIPRE.Tailored.Sofic.Measure.letterFG, MIPRE.Tailored.Sofic.Measure.wordFG,
  MIPRE.Tailored.Sofic.Measure.litB, MIPRE.Tailored.Sofic.Measure.passB,
  MIPRE.Tailored.Sofic.Measure.passSet, MIPRE.Tailored.Sofic.Measure.valGen,
  SubgroupTestValue.FiniteAction.finDesc, SubgroupTestValue.SubgroupTestData.valOn,
  SubgroupTestValue.SubgroupTestData.valErg}
\leanok
The value of a test $T$ against a probability measure $\mu$ on $\mathrm{Sub}(F)$ is the weighted
average over challenges of the measure of the set of subgroups passing the challenge. The
ergodic value $\mathrm{val}_{\mathrm{erg}}(T)$ is its supremum over the invariant random
subgroups. A finite action $\sigma$ gives the measure $\Phi(\sigma)$, the stabilizer of a uniform
point.
\end{definition}

\begin{lemma}[The value is a continuous functional; I:550]
\label{lem:val-on-continuous}
\uses{def:erg-value}
\lean{MIPRE.Tailored.Sofic.Measure.lift_wordFG, MIPRE.Tailored.Sofic.Measure.isClopen_passSet,
  SubgroupTestValue.SubgroupTestData.continuous_valOn, SubgroupTestValue.SubgroupTestData.valOn_of_eq,
  SubgroupTestValue.SubgroupTestData.valOn_finDesc, SubgroupTestValue.finDescIRS_subset_IRS}
\leanok
The value against a measure is weak-$*$ continuous, and the value against $\Phi(\sigma)$ is the
value of $\sigma$. Finitely described IRSs are IRSs.
\end{lemma}

\begin{proof}
\leanok
Each passing set is clopen in the product topology, so its measure is continuous in $\mu$; the
value is a finite weighted sum of these.
\end{proof}

\begin{corollary}[The sofic value is at most the ergodic value, with equality under
Aldous--Lyons; Corollary~I:550]
\label{cor:sof-eq-erg}
\uses{lem:val-on-continuous}
\lean{SubgroupTestValue.SubgroupTestData.valSof_le_valErg,
  SubgroupTestValue.SubgroupTestData.valSof_eq_valErg}
\leanok
For every test, $\mathrm{val}_{\mathrm{sof}}(T) \leq \mathrm{val}_{\mathrm{erg}}(T)$, and if
the Aldous--Lyons conjecture holds they are equal.
\end{corollary}

\begin{proof}
\leanok
The finitely described IRSs are IRSs; under the conjecture every IRS is a weak-$*$ limit of
finitely described ones, and the value is continuous.
\end{proof}

\begin{theorem}[Main Theorem~I (1): the sofic value is approximable from below; I:594]
\label{thm:main-theorem-i-lower}
\uses{lem:val-on-continuous}
\lean{MIPRE.Tailored.Sofic.Measure.Represents, MIPRE.Tailored.Sofic.Measure.ValidC,
  MIPRE.Tailored.Sofic.Measure.Represents.value_eq, MIPRE.Tailored.Sofic.Measure.decodeAct,
  MIPRE.Tailored.Sofic.Measure.codeOf, MIPRE.Tailored.Sofic.Measure.lowerSeq,
  MIPRE.Tailored.Sofic.Measure.dyadic, MIPRE.Tailored.Sofic.Measure.lowerSeq_le,
  MIPRE.Tailored.Sofic.Measure.lowerSeq_monotone, MIPRE.Tailored.Sofic.Measure.tendsto_lowerSeq,
  MIPRE.Tailored.Sofic.Measure.primrec_lowerSeq, MIPRE.Tailored.Sofic.Measure.mainTheoremI_one}
\leanok
There is a primitive recursive $\alpha(T, t)$ such that $\alpha(T,t)/2^t$ is non-decreasing in
$t$, at most $\mathrm{val}_{\mathrm{sof}}(T)$, and converges to it.
\end{theorem}

\begin{proof}
\leanok
Enumerate the finite actions by codes; $\alpha(T, t)/2^t$ is the largest value of the first $t$
codes' actions, rounded down to a multiple of $2^{-t}$.
\end{proof}

\begin{lemma}[From word indicators to invariant random subgroups]
\label{lem:irs-of-good}
\uses{def:erg-value,def:assoc-test}
\lean{MIPRE.Tailored.Sofic.Measure.valΩ, MIPRE.Tailored.Sofic.Measure.Good,
  MIPRE.Tailored.Sofic.Measure.wordEval, MIPRE.Tailored.Sofic.Measure.evW,
  MIPRE.Tailored.Sofic.Measure.conjW, MIPRE.Tailored.Sofic.Measure.toSub,
  MIPRE.Tailored.Sofic.Measure.irs_of_good, MIPRE.Tailored.Sofic.Measure.valOn_eq_valΩ}
\leanok
Index the indicator of a subgroup by words rather than group elements. A probability measure on
word indicators carried by the indicators that respect free reduction and the subgroup axioms,
and invariant under conjugation by each letter, is the image of an IRS with the same value
against every test.
\end{lemma}

\begin{proof}
\leanok
The good indicators are the complement of a countable union of clopen sets, and on them
evaluation at a reduced word is a subgroup indicator; push forward and check invariance on
generators.
\end{proof}

\begin{definition}[Pseudo-subgroup polytopes, discretized; I:927, I:1105]
\label{def:pseudo-subgroup-polytope}
\uses{def:subgroup-test,def:assoc-test,lem:irs-of-good}
\lean{MIPRE.Tailored.Sofic.Measure.wordsUpTo, MIPRE.Tailored.Sofic.Measure.stageW,
  MIPRE.Tailored.Sofic.Measure.stageD, MIPRE.Tailored.Sofic.Measure.goodP,
  MIPRE.Tailored.Sofic.Measure.allVecs, MIPRE.Tailored.Sofic.Measure.invErr,
  MIPRE.Tailored.Sofic.Measure.Feasible, MIPRE.Tailored.Sofic.Measure.numG,
  MIPRE.Tailored.Sofic.Measure.maxNum, MIPRE.Tailored.Sofic.Measure.Covers,
  MIPRE.Tailored.Sofic.Measure.upperSeq}
\leanok
At stage $t$ take the words of length at most $t+2$ in the first $n+t$ letters, and the grid
points of weight $M = P \cdot 2^t$ on its $P$ patterns. A grid point is feasible when every
pattern of positive weight is locally closed (contains the empty word, closed under cancelling a
pair, deleting a letter beyond the $n$ generators, and $u, v \mapsto u v^{-1}$ within the stage)
and conjugation by each generator moves it by at most $4P$ in $\ell^1$. The upper value at stage
$t$ is the best value of a feasible point, plus $2/2^t$, rounded up to a multiple of $2^{-t}$.
This replaces the paper's pseudo IRS polytope and its linear program by an exhaustive search over
a grid, with the invariance relaxed to absorb rounding.
\end{definition}

\begin{lemma}[The upper value bounds the ergodic value]
\label{lem:erg-le-upper}
\uses{def:pseudo-subgroup-polytope,lem:irs-of-good,def:assoc-test,thm:main-theorem-i-lower}
\lean{MIPRE.Tailored.Sofic.Measure.exists_rounding,
  MIPRE.Tailored.Sofic.Measure.exists_feasible_of_irs,
  MIPRE.Tailored.Sofic.Measure.valOn_le_upperSeq,
  MIPRE.Tailored.Sofic.Measure.valErg_le_upperSeq}
\leanok
For every stage $t$, $\mathrm{val}_{\mathrm{erg}}(T) \leq \mathrm{upper}(T, t)/2^t$.
\end{lemma}

\begin{proof}
\leanok
Round the restriction of an IRS to the stage's words to the grid: it is feasible, and its value
is within $2/2^t$ of the IRS's.
\end{proof}

\begin{lemma}[The upper value converges to the ergodic value]
\label{lem:upper-converge}
\uses{lem:erg-le-upper,lem:irs-of-good,def:assoc-test,lem:val-on-continuous,thm:main-theorem-i-lower}
\lean{MIPRE.Tailored.Sofic.Measure.isPiSystem_cylinders,
  MIPRE.Tailored.Sofic.Measure.eventually_valΩ_le, MIPRE.Tailored.Sofic.Measure.liftStage,
  MIPRE.Tailored.Sofic.Measure.valΩ_liftStage, MIPRE.Tailored.Sofic.Measure.tendsto_upperSeq}
\leanok
$\mathrm{upper}(T, t)/2^t \to \mathrm{val}_{\mathrm{erg}}(T)$.
\end{lemma}

\begin{proof}
\leanok
Lift the optimal feasible point of each stage to a measure on word indicators. The probability
measures on word indicators form a compact space, so a subsequence converges; the limit is
carried by the good indicators and is invariant, because the cylinder sets are a generating
$\pi$-system and the stage's conditions hold on ever more of them. By
Lemma~\ref{lem:irs-of-good} it is an IRS, and the value is continuous.
\end{proof}

\begin{theorem}[Main Theorem~I (2): the ergodic value is approximable from above; I:594]
\label{thm:main-theorem-i-upper}
\uses{lem:erg-le-upper,lem:upper-converge,def:assoc-test,lem:assoc-test-primrec,thm:main-theorem-i-lower}
\lean{MIPRE.Tailored.Sofic.Measure.primrec_upperSeq, MIPRE.Tailored.Sofic.Measure.upperMin,
  MIPRE.Tailored.Sofic.Measure.primrec_upperMin, MIPRE.Tailored.Sofic.Measure.valErg_le_upperMin,
  MIPRE.Tailored.Sofic.Measure.antitone_upperMin, MIPRE.Tailored.Sofic.Measure.mainTheoremI_two,
  MIPRE.Tailored.Sofic.Measure.ErgUpperApprox, MIPRE.Tailored.Sofic.Measure.MainTheoremITwo,
  MIPRE.Tailored.Sofic.Measure.ergUpperApprox, MIPRE.Tailored.Sofic.Measure.mainTheoremITwo}
\leanok
There is a primitive recursive $\beta(T, t)$ such that $\beta(T,t)/2^t$ is non-increasing in
$t$, at least $\mathrm{val}_{\mathrm{erg}}(T)$, and converges to it.
\end{theorem}

\begin{proof}
\leanok
The running minimum of the upper values, which is primitive recursive because the stage search
is.
\end{proof}

\begin{corollary}[Aldous--Lyons makes the sofic value approximable; Corollary~I:605]
\label{cor:al-sof-approximable}
\uses{cor:sof-eq-erg,thm:main-theorem-i-lower,thm:main-theorem-i-upper,def:sofic-approximable}
\lean{MIPRE.Tailored.Sofic.Measure.ceilDiv, MIPRE.Tailored.Sofic.Measure.approx_of_bounds,
  MIPRE.Tailored.Sofic.Measure.sofValueApproximable_of,
  MIPRE.Tailored.Sofic.Measure.sofValueApproximable_of_aldousLyons}
\leanok
If the Aldous--Lyons conjecture holds, the sofic value is approximable.
\end{corollary}

\begin{proof}
\leanok
Under the conjecture both sequences converge to the same value from either side; run them until
their gap is at most $1/(k+1)$, which happens, and output a point between them.
\end{proof}

\begin{theorem}[The Aldous--Lyons conjecture is false; Corollary~I:2144]
\label{thm:aldous-lyons-false}
\uses{def:aldous-lyons,thm:sofic-not-approximable,thm:main-theorem-ii,cor:al-sof-approximable}
\lean{MIPRE.Tailored.Sofic.Measure.aldous_lyons_false_of_upper,
  MIPRE.Tailored.Sofic.Measure.aldous_lyons_false_of, MIPRE.Tailored.Sofic.aldous_lyons_false}
\leanok
A halting reduction to tailored games refutes the Aldous--Lyons conjecture.
\end{theorem}

\begin{proof}
\leanok
Under the conjecture the sofic value is approximable (Corollary~\ref{cor:al-sof-approximable}),
against Theorem~\ref{thm:sofic-not-approximable} with Main Theorem~II.
\end{proof}

\begin{corollary}[The Aldous--Lyons conjecture is false, from tailored answer reduction]
\label{cor:aldous-lyons-false-ar}
\uses{thm:aldous-lyons-false,cor:tmip-re-from-ar}
\lean{MIPRE.Tailored.aldous_lyons_false_of_answerReduction}
\leanok
A tailored answer reduction at input level~$5$ refutes the Aldous--Lyons conjecture.
\end{corollary}

\begin{proof}
\leanok
Theorem~\ref{thm:aldous-lyons-false} with the halting reduction of
Corollary~\ref{cor:tmip-re-from-ar}.
\end{proof}

\begin{corollary}[The Aldous--Lyons conjecture is false]
\label{cor:aldous-lyons-false-final}
\uses{thm:aldous-lyons-false,thm:sofic-not-approximable,thm:main-theorem-ii,thm:tmip-re,cor:sofic-not-approximable-ar}
\lean{SubgroupTestValue.not_sofValueApproximable, SubgroupTestValue.aldous_lyons_false}
\leanok
The sofic value of subgroup tests is not approximable, and the Aldous--Lyons conjecture is false.
\end{corollary}

\begin{proof}
\leanok
Theorems~\ref{thm:sofic-not-approximable} and~\ref{thm:aldous-lyons-false} at the halting
reduction of Theorem~\ref{thm:tmip-re}.
\end{proof}
